\documentclass[8pt]{article}
\usepackage{graphicx} % Required for inserting images
\usepackage{amsthm}
\usepackage{amsmath}
\usepackage{amsfonts}
\usepackage{amsopn}
\usepackage{comment}
\usepackage{amssymb}
\usepackage{hyperref}
\usepackage{bussproofs}
\usepackage[all, 2cell]{xy}
\usepackage[all]{xy}
\usepackage{rotating}
\usepackage{lscape}
\usepackage{minted}
\usepackage{dsfont}
\usepackage{tikz}
\usepackage{tikz-cd}
\usepackage{stmaryrd}
\usepackage{multicol}
\usepackage{cmll}
\usepackage{wasysym}
\usetikzlibrary{shapes.geometric, arrows, positioning, backgrounds}
\usetikzlibrary{cd}
\makeatletter
\tikzcdset{
  eq node/.style={
    commutative diagrams/math mode=false, anchor=center},
  eq/.style={
    phantom,
    /tikz/every to/.append style={
      edge node={node[commutative diagrams/eq node]
        {\@eqnswtrue\make@display@tag\ltx@label{#1}}}}}}
\makeatother

\theoremstyle{definition}
\newtheorem{definition}{Definition}[section]

\theoremstyle{definition}
\newtheorem{theorem}{Theorem}[section]

\theoremstyle{definition}
\newtheorem{claim}{Claim}[section]

\theoremstyle{definition}
\newtheorem{ex}{Example}[section] 

\theoremstyle{definition}
\newtheorem{cons}{Construction}[section] 

\theoremstyle{definition}
\newtheorem{rem}{Remark}[section] 


\theoremstyle{definition}
\newtheorem{prop}{Proposition}[section]

\theoremstyle{definition}
\newtheorem{lemma}{Lemma}[section]

\theoremstyle{definition}
\newtheorem{fact}{Fact}[section]

\theoremstyle{definition}
\newtheorem{remark}{Remark}[section]

\theoremstyle{definition}
\newtheorem{notation}{Notation}[section]

\theoremstyle{definition}
\newtheorem{example}{Example}[section]

\theoremstyle{definition}
\newtheorem{exercise}{Exercise}[section]

\theoremstyle{definition}
\newtheorem{col}{Corollary}[section]

\theoremstyle{question}
\newtheorem{question}{Question}

\let\strokeL\L
\renewcommand\L{\mathbf{L}}

\DeclareMathOperator{\bang}{!}

\DeclareMathOperator*{\colim}{\operatorname{colim}}
\newcommand{\Ob}[1]{\operatorname{Ob}({\mathcal{#1}})}
\newcommand{\Cov}[1]{\operatorname{Cov}(#1)}

\newcommand{\circvee}{\mathbin{\text{\ooalign{\hfil$\vee$\hfil\cr\kern0.1ex$\bigcirc$}}}}
\newcommand{\circwedge}{\mathbin{\text{\ooalign{\hfil$\wedge$\hfil\cr\kern0.1ex$\bigcirc$}}}}

\title{Notes on Template Games and Differential Logic}
\author{Daniel Rogozin}
\date{ }

\newcommand{\IntHom}[1]{[#1]}

\begin{document}

\maketitle

\tableofcontents

\section{General Remarks}

The results from \cite{fiore2008cartesian} introduce the bicategorical interpretation of linear logic
based on distributors and generalised species. The idea is that every formula $A$ of linear logic is interpreted as
a small category $[\![A]\!]$ and every derivation tree
\begin{prooftree}
  \AxiomC{$\pi$}
  \noLine
  \UnaryInfC{$\:\:$}
  \noLine
  \UnaryInfC{$\ldots$}
  \noLine
  \UnaryInfC{$\:\:$}
  \UnaryInfC{$A_1, \ldots, A_n \vdash A$}
\end{prooftree}
is interpreted as a distributor (or a profunctor)
\[
[\![\pi]\!] : [\![A_1]\!] \otimes \ldots \otimes [\![A_n]\!] \not\to [\![A]\!].
\]
Recall that a distributor $M : A \not\to B$ is defined as a functor $M : A \times B^{op} \to {\bf Set}$.
The category ${\bf Dist}$ of small categories and distributors is symmetric monoidal with the tensor product defined
by the direct product of categories in ${\bf Cat}$, the category of small categories and functors, and the unit defined by the terminal 
(a single-element) category. ${\bf Dist}$ is also *-autonomous (and, moreover, compact closed)
with linear negation defined as $A \mapsto A^{op}$. We thus have the following isomorphism
\[
{\bf Dist}(A \otimes B, C) \cong {\bf Dist}(B, A^{op} \otimes C)
\]
natural $A$, $B$ and $C$, providing the currification in ${\bf Dist}$.

The exponential modality $A \mapsto \bang A$ is interpreted using the 2-monad ${\bf Sym} : {\bf Cat} \to {\bf Cat}$
sending each small category $\mathcal{C}$ to the free symmetric monoidal category generated by $\mathcal{C}$.
${\bf Sym}(\mathcal{C})$ has objects defined as finite sequences $A_1 \ldots A_n$ for objects from $\mathcal{C}$
for where $A_i \in \mathcal{C}$ for $1 \leq i \leq n$, and morphisms
\[
(\sigma, f_1 \dots f_n) : A_1 \ldots A_n \to B_1 \ldots B_n
\]
defined as pairs $(\sigma, f_1 \dots f_n)$ consists of a permutation $\sigma \in \mathcal{S}_n$ on the set
of $n$-elements along with a sequence of morphisms from $\mathcal{C}$ for $1 \leq k \leq n$
\[
f_k : A_k \to B_{\sigma(k)}.
\]
The category ${\bf Sym}(\mathcal{C})$, as strict monoidal, comes with a pair of functors 
$\otimes : {\bf Sym}(\mathcal{C}) \times {\bf Sym}(\mathcal{C}) \to {\bf Sym}(\mathcal{C})$ and $I_A : \mathds{1} \to {\bf Sym}(\mathcal{C})$
defined as the concatenation and the empty sequence respectively.
${\bf Sym}(\mathcal{C})$ with those two functors also forms a monoid in ${\bf Cat}$, and such a monoid
is commutative up to isomorphism $\gamma$ (the symmetry of the monoidal category):
\[
\begin{tikzcd}
  {\bf Sym}(\mathcal{C}) \times {\bf Sym}(\mathcal{C}) \ar[dr, "\otimes"', ""{name=otimes1}] \ar[rr, "(21)"] && {\bf Sym}(\mathcal{C}) \times {\bf Sym}(\mathcal{C}) \ar[dl, "\otimes", ""{name=otimes2}] \\
  & {\bf Sym}(\mathcal{C})
  \arrow[from=otimes1, to=otimes2, Rightarrow, "\alpha", shorten=6mm]
\end{tikzcd}
\]
where $(21)$ is the symmetry of the cartesian category ${\bf Sym}(\mathcal{C}) \times {\bf Sym}(\mathcal{C})$.
The exponential modality $A \mapsto \bang A$ is, thus, interpreted in ${\bf Dist}$ by turning ${\bf Sym}$ into a 2-comonad
\[
{\bf Sym} : {\bf Dist} \to {\bf Dist}
\]
using the distributivity law between ${\bf Sym}$ and the presheaf construction in ${\bf Cat}$. Further,
every functor $F : \mathcal{C} \to \mathcal{D}$ induces an adjoint pair $L_F \dashv R_F$ of distributors
$L_F : A \not\to B$ and $R_F : B \not\to A$ in the bicategory ${\bf Dist}$, where the distributors 
$L_F$ and $R_F$ are given as the following presheaves
\begin{center}
  $L_F(B, A) = \mathcal{D}(FB, A) : B^{op} \times A \to {\bf Set}$

  $R_F(A, B) = \mathcal{D}(A, FB) : A^{op} \times B \to {\bf Set}$
\end{center}

The comonoid structure $({\bf Sym}(\mathcal{C}), d_{\mathcal{C}}, e_{\mathcal{C}})$ of the exponential modality 
${\bf Sym}$ is given by the following distributors
\begin{center}
  $d_{\mathcal{C}} = R_{\otimes_{\mathcal{C}}} : {\bf Sym}(\mathcal{C}) \not\to {\bf Sym}(\mathcal{C}) \otimes {\bf Sym}(\mathcal{C})$

  $e_{\mathcal{C}} = R_{I_{\mathcal{C}}} : {\bf Sym}(\mathcal{C}) \not\to \mathds{1}$
\end{center}
right adjoint to $\otimes_{\mathcal{C}}$ and $I_{\mathcal{C}}$ respectively. Note that $({\bf Sym}(\mathcal{C}), d_{\mathcal{C}}, e_{\mathcal{C}})$
is not commutative, but only symmetric, so it is commutative only up to isomorphism:
\[
\begin{tikzcd}
& {\bf Sym}(\mathcal{C}) \ar[dl, "d_{\mathcal{C}}"', ""{name=d1}] \ar[dr, "d_{\mathcal{C}}", ""{name=d2}] \\
{\bf Sym}(\mathcal{C}) \otimes {\bf Sym}(\mathcal{C}) \ar[rr, "(21)"] && {\bf Sym}(\mathcal{C}) \otimes {\bf Sym}(\mathcal{C})
\arrow[from=d1, to=d2, Rightarrow, "\alpha", shorten=6mm]
\end{tikzcd}
\]
where $\gamma$ is the symmetry of ${\bf Sym}(\mathcal{C}) \otimes {\bf Sym}(\mathcal{C})$, which is a distributor in ${\bf Dist}$
as well.

The necessity to weaken the commutativity condition to an isomorphism is similar to the phenomenon of the Seely
isomorphism in linear logic:
\[
\bang (A \& B) \cong \bang A \otimes \bang B
\]
One can interpret this equation of linear logic as a distributor with the concatenation functor:
\[
\operatorname{concat}_{A,B} : {\bf Sym}(A) \times {\bf Sym}(B) \to {\bf Sym}(A + B) 
\]
that takes sequences $u = a_1 \ldots a_p$ and $v = b_1 \ldots b_q$ of objects from $A$ and $B$ respectively concatenates them
into the sequence of objects from $A + B$:
\[
\operatorname{concat}_{A,B}(u,v) = a_1 \ldots a_p \cdot b_1 \ldots b_q
\]
The distributor interpreting the Seely isomorphism
\[
{\bf Sym}(A \& B) \not\to {\bf Sym}(A \otimes B)
\]
is defined as the right adjoint of the concatenation functor where $A \& B$ stands for $A + B$
for the disjoint union of $A$ and $B$. Besides, $\operatorname{concat}_{A,B}$ is equivalence of categories injective
on objects, but it is not an isomorphism. Thus we replace here the Seely isomorphism with the Seely equivalence.

\section{On Models of Differential Linear Logic}

A good introduction to differential linear logic is \cite{ehrhard2018introduction}. If we interpret the
exponential modality as $A \mapsto {\bf Sym}(A)$, then one can see that the model of distributors is also a model
of differential linear logic since the exponential modality ${\bf Sym}(A)$ comes with the monoid structure in
${\bf Dist}$:
\begin{itemize}
  \item $m_A = L_{\otimes_A} : {\bf Sym}(A) \otimes {\bf Sym}(A) \not\to {\bf Sym}(A)$,
  \item $u_A = L_{I_A} : \mathds{1} \not\to {\bf Sym}(A)$.
\end{itemize}
whose multiplication and unit are the left adjoint distributors associated to the monoid structure of ${\bf Sym}(A)$ in ${\bf Cat}$.
Note that the monoid $({\bf Sym}(A), m_A, u_A)$ and the comonoid $({\bf Sym}(A), d_A, e_A)$ are the left and right adjoint ``avatars'' of the 
monoid structure $({\bf Sym}(A), \otimes_A, I_A)$ in ${\bf Cat}$. The monoid and the comonoid structures on ${\bf Sym}(A)$ define a bialgebra, but the bialgebra
structure is only up to an invertible natural transformation:
\[
\begin{tikzcd}
  && |[alias=symA]| {\bf Sym}(A) \ar[drr, "d_A"] \\
{\bf Sym}(A)^{\otimes 2} \ar[dr, "d_A \otimes d_A"'] \ar[urr, "m_A"] &&&& {\bf Sym}(A)^{\otimes 2} \\
& {\bf Sym}(A)^{\otimes 4} \ar[rr, "(1324)"', ""{name=1324}] && {\bf Sym}(A)^{\otimes 4} \ar[ur, "m_A \otimes m_A"']
\arrow[from=symA,to=1324,Rightarrow,"\sim"]
\end{tikzcd}
\]
The explanation is that the disjoint sum $A + B$ of small categories $A$ and $B$ coincide with their cartesian sum $A \oplus B$
and their cartesian product $A \& B$. Therefore, every object $A$ is equipped with a monoid and a comonoid structure:
\begin{itemize}
  \item $\nabla_A : A \oplus A \not\to A$,
  \item $\nabla^{0}_A : 0 \not\to A$,
  \item $\Delta_A : A \not\to A \oplus A$,
  \item $\Delta^{0}_A :  A \not\to 0$.
\end{itemize}
defined as the left and right adjoint distributors to the canonical functors $A + A \to A$ and $0 \to A$,
where $0$ is the initial category in ${\bf Cat}$. The multiplication $\nabla_A$ and the comultiplication $\Delta_A$ are equipped with a bialgebra
structure in ${\bf Dist}$, up to an invertible natural transformation:
\[
\begin{tikzcd}
  && |[alias=A]| A \ar[drr, "\Delta_A"] \\
A^{\oplus 2} \ar[dr, "\Delta_A \oplus \Delta_A"'] \ar[urr, "\nabla_A"] &&&& A^{\oplus 2} \\
& A^{\oplus 4} \ar[rr, "(1324)"', ""{name=1324}] && A^{\oplus 4} \ar[ur, "\nabla_A \oplus \nabla_A"']
\arrow[from=A, to=1324, Rightarrow, "\sim"]
\end{tikzcd}
\]

The exponential modality $A \mapsto {\bf Sym}(A)$ along the family of Seely equivalences
\[
{\bf Sym}(A^{\oplus n}) \not\to ({\bf Sym}(A))^{\otimes n}
\]
defines a lax and oplax symmetric monoidal bifunctor
\[
{\bf Sym} : ({\bf Dist}, \oplus, 0) \to ({\bf Dist}, \otimes, \mathds{1})
\]

The differential in DiLL is interpreted the following way in the category of distributors. One starts with a family of functors
\[
\eta_A : A \to {\bf Sym}(A)
\]
defining the unit of the monad $A \mapsto {\bf Sym}(A)$. One can postcompose the functor with $\otimes_A$ to obtain the functor:
\[
\operatorname{append}_A : (a, v) \mapsto a \cdot v \in {\bf Sym}(A) : A \times {\bf Sym}(A) \to {\bf Sym}(A)
\]
The differential is interpreted in ${\bf Dist}$ as the left adjoint distributor
\[
\partial_A : A \otimes {\bf Sym}(A) \not\to {\bf Sym}(A)
\]
associated to the $\operatorname{append}_A$, whereas the dereliction and the codereliction morphisms:
\begin{itemize}
  \item $\operatorname{der}_A : {\bf Sym}(A) \not\to A$,
  \item $\operatorname{coder}_A : A \not\to {\bf Sym}(A)$
\end{itemize}
are the left and right adjoint distributors associated to $\eta_A$.

The composite of the codereliction with the comultiplication
\[
A \xrightarrow{\operatorname{coder}_A } {\bf Sym}(A) \xrightarrow{d_A} {\bf Sym}(A) \otimes {\bf Sym}(A)
\]
is equal to the disjoint union of the following two composites:
\begin{center}
  $A \xrightarrow{\operatorname{coder}_A } {\bf Sym}(A) \xrightarrow{u_A \otimes 1_{\bf Sym}(A)} {\bf Sym}(A) \otimes {\bf Sym}(A)$

  $A \xrightarrow{\operatorname{coder}_A } {\bf Sym}(A) \xrightarrow{1_{\bf Sym}(A) \otimes u_A}{\bf Sym}(A) \otimes {\bf Sym}(A)$
\end{center}
where the disjoint product of distributors $M, N : A \not\to B$ is the distributor $M + N$ defined using the convolution product
associated to the additive comonoid and monoid structures on $A$ and $B$:
\[
A \xrightarrow{\Delta_A} A \oplus A \xrightarrow{M + N} B \oplus B \xrightarrow{\nabla_B} B
\]
or defined as the presheaf:
\[
M + N : (b,a) \mapsto M(b,a) \sqcup N(b,a)
\]

The property of the dereliction ensures that the Leibniz rule is satisfied, in the technical sense that the composite distributor
\[
A \otimes {\bf Sym}(A) \xrightarrow{\partial_A} {\bf Sym}(A) \xrightarrow{d_A} {\bf Sym}(A) \otimes {\bf Sym}(A)
\]
is isomorphic to the sum of the two distributors obtained by precomposing the distributor:
\[
A \otimes {\bf Sym}(A) \xrightarrow{1_A \otimes d_A} A \otimes {\bf Sym}(A) \otimes {\bf Sym}(A)
\]
with the two distributors represented below:
\[
\begin{tikzcd}
  A \otimes {\bf Sym}(A) \otimes {\bf Sym}(A) \ar[rr, "\partial_A \otimes 1_{{\bf Sym}(A)}"] \ar[dr, "\operatorname{symmetry}"'] && {\bf Sym}(A) \otimes {\bf Sym}(A) \\
  & {\bf Sym}(A) \otimes A \otimes {\bf Sym}(A) \ar[ur, "1_{{\bf Sym}(A)} \otimes \partial_A"']
\end{tikzcd}
\]

The general idea is to refine the exponential modality $A \mapsto {\bf Sym}(A)$ of the models of distributors in order to obtain
the game semantic interpretation of differential linear logic. The difference between two models (games and distributors) is that
template games are based on functorial spans between categories instead of distributors. This would allow providing bridges with the canonical Quillen
model structure on the category ${\bf Cat}$.

The concept of a template game was introduced as a unified framework to construct various *-autonomous bicategories ${\bf Games}(\davidsstar)$.
Each such ${\bf Games}(\davidsstar)$ is built in a unified manner using a \emph{synchronisation template}. Such a template is denoted as 
$\davidsstar$ and called an \emph{anchor}. The template $\davidsstar$ expresses the scheduling policy of a specific regime of games
and strategies. For example:
\begin{itemize}
  \item $\davidsstar_{alt}$ for sequential alternating games, 
  \item $\davidsstar_{conc}$ for concurrent non-alternating games,
  \item $\davidsstar_{span}$ for functorial spans without scheduling.
\end{itemize}
The higher algebraic structure of the bicategory ${\bf Games}(\davidsstar)$ reflects the simpler combinatorial structure of the template $\davidsstar$.
The construction of ${\bf Games}(\davidsstar)$ based on the hypothesis that the template $\davidsstar$ defines an internal category in 
a category $\mathbb{S}$ with finite limits, and, most of the time, $\mathbb{S} = {\bf Cat}$. Further, we will be calling internal category
(functor, natural transformation) and $\mathbb{S}$-category (functor, natural transformation). ${\bf Cat}(\mathbb{S})$ stands for the 2-category 
of internal two categories.

\begin{theorem}
  The bicategory ${\bf Games}(\davidsstar)$ of games, strategies and simulations is *-autonomous whenever the $\mathbb{S}$-category is span-monoidal
  *-autonomous.
\end{theorem}

To obtain a template game model of differential linear logic with the exponential modality $A \mapsto \bang A$, one needs to describe what the extra-structure of the template $\davidsstar$ looks like.

\bibliographystyle{alpha}
\bibliography{Text}

\end{document}